%----------------------------------------------------------------------------------------
%  MyST LaTeX template for the QM lecture notes
%  Style adapted from Paul Skrzypczyk's QM1 notes (2022), themselves based on
%  "The Legrand Orange Book" by Mathias Legrand / Vel (CC BY-NC-SA 3.0, LaTeXTemplates.com)
%----------------------------------------------------------------------------------------
\documentclass[11pt,oneside,openany]{book}

%----------------------------------------------------------------------------------------
%	PACKAGES
%----------------------------------------------------------------------------------------
\usepackage[utf8]{inputenc}
\usepackage[T1]{fontenc}
\usepackage{graphicx}
\usepackage{tikz}
\usetikzlibrary{shapes.geometric}
\usepackage[shortlabels]{enumitem}
\setlist{nolistsep}
\usepackage{booktabs}
\usepackage{xcolor}
\usepackage{xstring}
\usepackage{amsmath,amsfonts,amssymb,amsthm}
\usepackage{microtype}

% Dirac notation as used on the website (KaTeX provides these natively)
\providecommand{\ket}[1]{\lvert #1 \rangle}
\providecommand{\bra}[1]{\langle #1 \rvert}
\providecommand{\braket}[2]{\langle #1 \vert #2 \rangle}

%----------------------------------------------------------------------------------------
%	COLOURS
%----------------------------------------------------------------------------------------
\definecolor{brightblue}{HTML}{008DCC}
\definecolor{brightred}{HTML}{C9002F}
\definecolor{UoBred}{HTML}{B01C2E}
\definecolor{UoBgrey}{HTML}{E5E6E5}
\definecolor{UoBorange}{HTML}{EE7219}
\definecolor{UoBbrightblue}{HTML}{0CC6DE}
\definecolor{exuorange}{HTML}{FF582D}
\definecolor{webblue}{HTML}{0076df}
\definecolor{maincolor}{HTML}{1D4ED8}   % a shade deeper than the figures' wavefunction blue (#2563EB); chosen 9 Nov 2026

\newcommand\alert[1]{\textcolor{maincolor}{\textbf{#1}}}
\newcommand\lesson[1]{\textcolor{maincolor}{\textit{#1}}}

%----------------------------------------------------------------------------------------
%	MARGINS
%----------------------------------------------------------------------------------------
\usepackage{geometry}
\geometry{
	paper=a4paper,
	top=2.5cm,
	bottom=2.5cm,
	left=2.5cm,
	right=2.5cm,
	headheight=14pt,
	footskip=0cm,
	headsep=10pt,
}

%----------------------------------------------------------------------------------------
%	FONTS: Times for the body, Avant Garde for headings
%----------------------------------------------------------------------------------------
\usepackage{times}
\usepackage{FiraSans}   % headings and labels (texlive-fonts-extra); the body stays Times

%----------------------------------------------------------------------------------------
%	MAIN TABLE OF CONTENTS
%----------------------------------------------------------------------------------------
\usepackage{titletoc}
\contentsmargin{0cm}

\titlecontents{part}
	[0cm]
	{\addvspace{20pt}\bfseries}
	{}
	{}
	{}

\titlecontents{chapter}
	[1.25cm]
	{\addvspace{12pt}\large\sffamily\bfseries}
	{\color{maincolor}\contentslabel[\Large\thecontentslabel]{1.25cm}\color{maincolor}}
	{\color{maincolor}}
	{\color{maincolor}\normalsize\;\titlerule*[.5pc]{.}\;\thecontentspage}

\titlecontents{section}
	[1.25cm]
	{\addvspace{3pt}\sffamily\bfseries}
	{\contentslabel[\thecontentslabel]{1.25cm}}
	{}
	{\hfill\color{black}\thecontentspage}

\titlecontents{subsection}
	[1.25cm]
	{\addvspace{1pt}\sffamily\small}
	{\contentslabel[\thecontentslabel]{1.25cm}}
	{}
	{\ \titlerule*[.5pc]{.}\;\thecontentspage}

\makeatletter
\newcommand*{\toccontents}{\@starttoc{toc}}
\makeatother

%----------------------------------------------------------------------------------------
%	HEADERS AND FOOTERS
%----------------------------------------------------------------------------------------
\usepackage{fancyhdr}
\pagestyle{fancy}
\renewcommand{\chaptermark}[1]{\markboth{\sffamily\normalsize\bfseries\color{maincolor}\chaptername\ \thechapter.\color{black}\ #1}{}}
\renewcommand{\sectionmark}[1]{\markright{\sffamily\normalsize\color{maincolor}\thesection\color{black}\hspace{5pt}#1}{}}
\fancyhf{}
\fancyhead[R]{\sffamily\normalsize\thepage}
\fancyhead[L]{\leftmark}
\renewcommand{\headrulewidth}{0.5pt}
% every page carries the licence, so that a page on its own still says what it is and whose it is
\newcommand{\licenceline}{\sffamily\scriptsize\color{gray}[-doc.title-] [-options.subtitle_text-] \textperiodcentered\ \textcopyright\ [# for author in doc.authors #][-author.name-][# if not loop.last #], [# endif #][# endfor #] \textperiodcentered\ CC BY-NC-ND 4.0 \textperiodcentered\ \url{https://paulskrzypczyk.github.io/qm-lecturenotes-new/}}
\fancyfoot[C]{\licenceline}
\fancypagestyle{plain}{\fancyhead{}\fancyfoot[C]{\licenceline}\renewcommand{\headrulewidth}{0pt}}

%----------------------------------------------------------------------------------------
%	THEOREM STYLES
%----------------------------------------------------------------------------------------
\makeatletter
\newtheoremstyle{maincolornumbox}
{0pt}{0pt}{\normalfont}{}
{\small\bf\sffamily\color{maincolor}}{\;}{0.25em}
{\small\sffamily\color{maincolor}\thmname{#1}\nobreakspace\thmnumber{\@ifnotempty{#1}{}\@upn{#2}}%
\thmnote{\nobreakspace\the\thm@notefont\sffamily\bfseries\color{black}---\nobreakspace#3.}}

\newtheoremstyle{blacknumbox}
{0pt}{0pt}{\normalfont}{}
{\small\bf\sffamily}{\;}{0.25em}
{\small\sffamily\thmname{#1}\nobreakspace\thmnumber{\@ifnotempty{#1}{}\@upn{#2}}%
\thmnote{\nobreakspace\the\thm@notefont\sffamily\bfseries---\nobreakspace#3.}}
\makeatother

\theoremstyle{blacknumbox}
\newtheorem{definitionT}{Definition}[chapter]
\newtheorem*{eqT*}{}

%----------------------------------------------------------------------------------------
%	COLOURED BOXES
%----------------------------------------------------------------------------------------
\usepackage[framemethod=default]{mdframed}

% Definition / key-statement box: thick blue bar on the left
\newmdenv[skipabove=7pt,skipbelow=7pt,
rightline=false,leftline=true,topline=false,bottomline=false,
linecolor=maincolor,innerleftmargin=5pt,innerrightmargin=5pt,
innertopmargin=0pt,innerbottommargin=0pt,leftmargin=0cm,rightmargin=0cm,
linewidth=4pt]{dBox}

% Equation box: grey background
\newmdenv[skipabove=7pt,skipbelow=7pt,
rightline=false,leftline=false,topline=false,bottomline=false,
backgroundcolor=UoBgrey!50,linecolor=maincolor,
innerleftmargin=5pt,innerrightmargin=5pt,innertopmargin=5pt,innerbottommargin=5pt,
leftmargin=0cm,rightmargin=0cm,linewidth=4pt]{eqBox}

\newenvironment{definition}{\begin{dBox}\begin{definitionT}}{\end{definitionT}\end{dBox}}
\newenvironment{boxedeq}{\begin{eqBox}\begin{eqT*}}{\end{eqT*}\end{eqBox}}

% Blockquotes (key statements in the notes) use the blue-barred box
\renewenvironment{quote}{\begin{dBox}\itshape}{\end{dBox}}

%----------------------------------------------------------------------------------------
%	REMARK ENVIRONMENT  (used for MyST asides)
%----------------------------------------------------------------------------------------
\newenvironment{remark}{\par\vspace{10pt}\small
\begin{list}{}{
\leftmargin=35pt
\rightmargin=25pt}\item\ignorespaces
\makebox[{-2.5pt}]{\begin{tikzpicture}[overlay]
\node[draw=maincolor,line width=1pt,circle,fill=UoBgrey!50,font=\sffamily\bfseries,inner sep=2pt,outer sep=0pt] at (-15pt,0pt){!};\end{tikzpicture}}
\advance\baselineskip -1pt \itshape}{\end{list}\vskip5pt}

%----------------------------------------------------------------------------------------
%	SECTION NUMBERING IN THE MARGIN
%----------------------------------------------------------------------------------------
\makeatletter
\renewcommand{\@seccntformat}[1]{\llap{\textcolor{maincolor}{\csname the#1\endcsname}\hspace{1em}}}
\renewcommand{\section}{\@startsection{section}{1}{\z@}
{-4ex \@plus -1ex \@minus -.4ex}
{1ex \@plus.2ex }
{\normalfont\large\sffamily\bfseries}}
\renewcommand{\subsection}{\@startsection {subsection}{2}{\z@}
{-3ex \@plus -0.1ex \@minus -.4ex}
{0.5ex \@plus.2ex }
{\normalfont\sffamily\bfseries}}
\renewcommand{\subsubsection}{\@startsection {subsubsection}{3}{\z@}
{-2ex \@plus -0.1ex \@minus -.2ex}
{.2ex \@plus.2ex }
{\normalfont\small\sffamily\bfseries}}
\renewcommand\paragraph{\@startsection{paragraph}{4}{\z@}
{-2ex \@plus-.2ex \@minus .2ex}
{.1ex}
{\normalfont\small\sffamily\bfseries}}

%----------------------------------------------------------------------------------------
%	CHAPTER HEADINGS
%----------------------------------------------------------------------------------------
\newcommand{\autodot}{.}
\newlength\chaptertitleheight
\newsavebox\chaptertitlebox

\def\@makechapterhead#1{%
{\parindent \z@ \raggedright \normalfont
\ifnum \c@secnumdepth >\m@ne
\setbox\chaptertitlebox=\hbox{%
\parbox{15cm}{\huge\sffamily\bfseries\color{black}\thechapter\autodot~#1\strut}}
\setlength\chaptertitleheight{\dimexpr\ht\chaptertitlebox+\dp\chaptertitlebox}
\if@mainmatter
\begin{tikzpicture}[remember picture,overlay]
	\node at (current page.north west){
		\begin{tikzpicture}[remember picture,overlay]
			\draw[anchor=west] (\Gm@lmargin+.3cm,-3cm) node [text width=15cm] {
			\huge\sffamily\bfseries\color{maincolor}\thechapter\autodot~\color{black}#1\strut};
			\draw[anchor=west,line width=2pt,maincolor]
				(\Gm@lmargin+.3cm,-3cm-\chaptertitleheight/2) -- (\Gm@lmargin+22cm,-3cm-\chaptertitleheight/2);
		\end{tikzpicture}};
\end{tikzpicture}
\else
\begin{tikzpicture}[remember picture,overlay]
	\node at (current page.north west){
	\begin{tikzpicture}[remember picture,overlay]
		\draw[anchor=west] (\Gm@lmargin+.3cm,-9cm) node [text width=15cm] {
		\huge\sffamily\bfseries\color{black}#1\strut};
		\draw[anchor=west,line width=2pt,maincolor]
				(\Gm@lmargin+.3cm,-3cm-\chaptertitleheight/2) -- (\Gm@lmargin+22cm,-3.5cm);
	\end{tikzpicture}};
	\end{tikzpicture}
\fi\fi\par\vspace*{50\p@}}}
% unnumbered chapters (the Bibliography) get the same banner, without a number
\def\@makeschapterhead#1{%
{\parindent \z@ \raggedright \normalfont
\setbox\chaptertitlebox=\hbox{\parbox{15cm}{\huge\sffamily\bfseries\color{black}#1\strut}}
\setlength\chaptertitleheight{\dimexpr\ht\chaptertitlebox+\dp\chaptertitlebox}
\begin{tikzpicture}[remember picture,overlay]
	\node at (current page.north west){
		\begin{tikzpicture}[remember picture,overlay]
			\draw[anchor=west] (\Gm@lmargin+.3cm,-3cm) node [text width=15cm] {
			\huge\sffamily\bfseries\color{black}#1\strut};
			\draw[anchor=west,line width=2pt,maincolor]
				(\Gm@lmargin+.3cm,-3cm-\chaptertitleheight/2) -- (\Gm@lmargin+22cm,-3cm-\chaptertitleheight/2);
		\end{tikzpicture}};
\end{tikzpicture}
\par\vspace*{50\p@}}}
\makeatother

% MyST passes the page titles "Chapter N: Title" / "Appendix N: Title".
% Strip the prefix (LaTeX numbers the chapters itself) and switch to
% appendix numbering when the first appendix is reached.
\newif\ifqmappendix
\let\qmorigchapter\chapter
\RenewDocumentCommand{\chapter}{s o m}{%
  \IfBooleanTF{#1}{\qmorigchapter*{#3}}{%
    \IfBeginWith{#3}{Chapter }{%
      \StrBehind{#3}{: }[\qmtitle]%
      \expandafter\qmorigchapter\expandafter{\qmtitle}%
    }{%
      \IfBeginWith{#3}{Appendix }{%
        \ifqmappendix\else\appendix\qmappendixtrue\fi%
        \StrBehind{#3}{: }[\qmtitle]%
        \expandafter\qmorigchapter\expandafter{\qmtitle}%
      }{%
        \qmorigchapter{#3}%
      }%
    }%
  }%
}

%----------------------------------------------------------------------------------------
%	BOLD TEXT IS BLUE (the \alert of the old notes)
%----------------------------------------------------------------------------------------
\let\qmtextbf\textbf
\renewcommand{\textbf}[1]{\qmtextbf{\textcolor{maincolor}{#1}}}

%----------------------------------------------------------------------------------------
%	FIGURES
%----------------------------------------------------------------------------------------
\usepackage{flafter}            % never place a float before the text that introduces it
\setcounter{topnumber}{3}
\renewcommand{\topfraction}{0.9}
\renewcommand{\textfraction}{0.05}
\renewcommand{\arraystretch}{1.5} % breathing room in tables (they contain fractions)

%----------------------------------------------------------------------------------------
%	DATE FORMAT ON THE TITLE PAGE (as in the old notes: "Thursday 7th September, 2023")
%----------------------------------------------------------------------------------------
\usepackage[UKenglish]{datetime}
\newdateformat{versiondate}{\dayofweekname{\THEDAY}{\THEMONTH}{\THEYEAR} \ordinaldate{\THEDAY} \monthname[\THEMONTH], \THEYEAR}

%----------------------------------------------------------------------------------------
%	MyST IMPORTS (packages it needs, and the math macros from myst.yml)
%----------------------------------------------------------------------------------------
[- IMPORTS -]

%----------------------------------------------------------------------------------------
%	LINKS
%----------------------------------------------------------------------------------------
\usepackage{natbib}   % author-year citations of the original papers (references.bib)
\usepackage{hyperref}
\hypersetup{hidelinks,colorlinks=true,breaklinks=true,urlcolor=maincolor,bookmarksopen=false,linkcolor=exuorange,
  pdftitle={[-doc.title-] Lecture Notes},pdfauthor={[# for author in doc.authors #][-author.name-][# if not loop.last #], [# endif #][# endfor #]}}

\begin{document}

%----------------------------------------------------------------------------------------
%	TITLE PAGE
%----------------------------------------------------------------------------------------
\begingroup
\thispagestyle{empty}
\enlargethispage{-3.4cm}   % keep the contents list clear of the licence and logo at the foot of the title page
\begin{tikzpicture}[remember picture,overlay]
\draw[anchor=north west] (current page.north) node [fill=UoBgrey,text opacity=1,inner sep=1cm,trapezium,trapezium angle=-13+90] (p1) {
	\Huge\bfseries\sffamily\parbox[c][][t]{0.85\paperwidth}{{\color{maincolor}\hphantom{.} \hspace{2cm} [-doc.title-]
	\\
	[10pt]{\huge \hphantom{.} \hspace{2.1cm}[-options.subtitle_text-]}}
\\[10pt]{\Large \hphantom{.} \hspace{2cm}[# for author in doc.authors #][-author.name-][# if not loop.last #], [# endif #][# endfor #] \hspace{5.2cm} [-options.year-]{\small\mdseries \\ \hphantom{.} \hspace{1.8cm} version: \versiondate\today}}}};
\draw [line width=0.75mm, maincolor ](p1.south west) -- ++(-12,0);
\draw [line width=0.75mm, maincolor ](p1.south west) -- ++(12,0);
% the University's logo, bottom right of the title page, beside the licence
\node[anchor=south east, inner sep=0pt] at ([xshift=-2.2cm, yshift=1.7cm]current page.south east) {\includegraphics[width=4.6cm]{UoB_RGB_24.png}};
% the licence, bottom left of the title page
\node[anchor=south west, inner sep=0pt, text width=0.55\paperwidth] at ([xshift=2.2cm, yshift=1.6cm]current page.south west) {\sffamily\small\color{gray}%
  \textcopyright\ [# for author in doc.authors #][-author.name-][# if not loop.last #], [# endif #][# endfor #], School of Physics, University of Bristol. These notes are released under the Creative Commons Attribution-NonCommercial-NoDerivatives 4.0 International licence (CC BY-NC-ND 4.0): see the final page. The current version is always at \url{https://paulskrzypczyk.github.io/qm-lecturenotes-new/}.};
\end{tikzpicture}
\vspace*{120pt}

{
\hypersetup{hidelinks,colorlinks=false,breaklinks=true,urlcolor=maincolor,bookmarksopen=false,linkcolor=maincolor}
\toccontents
}

\vfill
\endgroup

\pagestyle{fancy}

%----------------------------------------------------------------------------------------
%	CONTENT
%----------------------------------------------------------------------------------------
[- CONTENT -]

%----------------------------------------------------------------------------------------
%	BIBLIOGRAPHY (the original papers cited in the notes; MyST copies references.bib to main.bib)
%----------------------------------------------------------------------------------------
[# if doc.bibliography #]
\cleardoublepage
\phantomsection
\addcontentsline{toc}{chapter}{Bibliography}
\chapter*{Bibliography}
\markboth{Bibliography}{Bibliography}
\bibliographystyle{plainnat}
\renewcommand{\bibsection}{}
\bibliography{main}
[# endif #]

%----------------------------------------------------------------------------------------
%	LICENCE (the last page; the full legal code is on the website's licence page)
%----------------------------------------------------------------------------------------
\cleardoublepage
\phantomsection
\addcontentsline{toc}{chapter}{Licence}
\chapter*{Licence}
\markboth{Licence}{Licence}
\noindent These lecture notes, \emph{[-doc.title-] [-options.subtitle_text-]}, are \textcopyright\ [# for author in doc.authors #][-author.name-][# if not loop.last #], [# endif #][# endfor #], School of Physics, University of Bristol, and are released under the \textbf{Creative Commons Attribution-NonCommercial-NoDerivatives 4.0 International} licence (CC BY-NC-ND 4.0).

\medskip\noindent Under this licence you are free to \textbf{share} the notes: to copy and redistribute them in any medium or format, for any non-commercial purpose, including in a teaching setting. In return you must keep to the following terms.
\begin{description}[leftmargin=4.3cm, style=sameline]
  \item[Attribution] You must give appropriate credit (the author, the title and the address below), provide a link to the licence, and indicate if any changes were made. You may do so in any reasonable manner, but not in any way that suggests the author endorses you or your use.
  \item[NonCommercial] You may not use the material for commercial purposes.
  \item[NoDerivatives] If you remix, transform or build upon the material, you may not distribute the modified material.
  \item[No further restrictions] You may not apply legal terms or technological measures that legally restrict others from doing anything the licence permits.
\end{description}
\noindent The licence does not restrict uses that are permitted anyway by an exception or limitation to copyright, such as fair dealing for private study or criticism, and it gives no warranties. The full legal text of the licence is at \url{https://creativecommons.org/licenses/by-nc-nd/4.0/legalcode}, and is reproduced on the licence page of the online notes.

\medskip\noindent The notes are maintained online at \url{https://paulskrzypczyk.github.io/qm-lecturenotes-new/}, where the figures are interactive and where corrections and comments can be submitted; this PDF is a snapshot of that site, dated on the title page. The University of Bristol logo is a trade mark of the University of Bristol and is not covered by the licence.

\medskip\noindent\textbf{Suggested citation.} [# for author in doc.authors #][-author.name-][# if not loop.last #], [# endif #][# endfor #], \emph{[-doc.title-]: [-options.subtitle_text-]}, University of Bristol, [-options.year-], \url{https://paulskrzypczyk.github.io/qm-lecturenotes-new/}.

\end{document}
